% ============================================================
% Question Bank: ch07-01 Parts of a Circle
% Topic Title: Parts of a Circle
% CCSS: 7.G.B.4
% Questions: 27 (15 MC + 6 SA + 6 Graphical)
% ============================================================

% --- Multiple Choice Questions (15) ---

\begin{practiceQuestion}{7-1-q01}{mc}
\begin{questionText}
What is a line segment that has both endpoints on a circle called?
\end{questionText}
\choiceA{Radius}
\choiceB{Diameter}
\choiceC{Chord}
\choiceD{Center}
\correctAnswer{C}
\explanation{A chord is a line segment whose two endpoints both lie on the circle. A diameter is a special chord that passes through the center.}
\end{practiceQuestion}

\begin{practiceQuestion}{7-1-q02}{mc}
\begin{questionText}
A circular trampoline has a radius of $9$ feet. What is its diameter?
\end{questionText}
\choiceA{$4.5$ feet}
\choiceB{$9$ feet}
\choiceC{$18$ feet}
\choiceD{$27$ feet}
\correctAnswer{C}
\explanation{The diameter is twice the radius: $d = 2r = 2 \times 9 = 18$ feet.}
\end{practiceQuestion}

\begin{practiceQuestion}{7-1-q03}{mc}
\begin{questionText}
A manhole cover has a diameter of $22$ inches. What is its radius?
\end{questionText}
\choiceA{$44$ inches}
\choiceB{$22$ inches}
\choiceC{$11$ inches}
\choiceD{$5.5$ inches}
\correctAnswer{C}
\explanation{The radius is half the diameter: $r = d \div 2 = 22 \div 2 = 11$ inches.}
\end{practiceQuestion}

\begin{practiceQuestion}{7-1-q04}{mc}
\begin{questionText}
Which best describes the center of a circle?
\end{questionText}
\choiceA{The point equidistant from every point on the circle}
\choiceB{The longest chord of the circle}
\choiceC{The distance around the circle}
\choiceD{A line segment from one side to the other}
\correctAnswer{A}
\explanation{The center of a circle is the point that is the same distance from every point on the circle. That constant distance is the radius.}
\end{practiceQuestion}

\begin{practiceQuestion}{7-1-q05}{mc}
\begin{questionText}
A bicycle wheel has a diameter of $26$ inches. Ravi says the radius is $13$ inches. Is he correct?
\end{questionText}
\choiceA{No, the radius is $52$ inches.}
\choiceB{No, the radius is $26$ inches.}
\choiceC{Yes, because radius $= d \div 2$.}
\choiceD{Yes, because radius $= d \times 2$.}
\correctAnswer{C}
\explanation{The radius is half the diameter: $r = 26 \div 2 = 13$ inches. Ravi is correct because the formula is $r = d \div 2$.}
\end{practiceQuestion}

\begin{practiceQuestion}{7-1-q06}{mc}
\begin{questionText}
A pizza pan has a radius of $5.5$ inches. What is the diameter of the pan?
\end{questionText}
\choiceA{$2.75$ inches}
\choiceB{$5.5$ inches}
\choiceC{$8.5$ inches}
\choiceD{$11$ inches}
\correctAnswer{D}
\explanation{The diameter is twice the radius: $d = 2r = 2 \times 5.5 = 11$ inches.}
\end{practiceQuestion}

\begin{practiceQuestion}{7-1-q07}{mc}
\begin{questionText}
Which statement about all radii of the same circle is always true?
\end{questionText}
\choiceA{They are parallel to each other.}
\choiceB{They are all equal in length.}
\choiceC{They are all perpendicular to the diameter.}
\choiceD{They are each half the circumference.}
\correctAnswer{B}
\explanation{Every radius of a given circle extends from the center to a point on the circle, so all radii of the same circle have the same length.}
\end{practiceQuestion}

\begin{practiceQuestion}{7-1-q08}{mc}
\begin{questionText}
A circular clock face has a radius of $15$ cm. An ant walks straight across the clock through the center. How far does the ant walk?
\end{questionText}
\choiceA{$7.5$ cm}
\choiceB{$15$ cm}
\choiceC{$30$ cm}
\choiceD{$45$ cm}
\correctAnswer{C}
\explanation{Walking across the clock through the center covers the diameter. The diameter is $d = 2r = 2 \times 15 = 30$ cm.}
\end{practiceQuestion}

\begin{practiceQuestion}{7-1-q09}{mc}
\begin{questionText}
A circle has a diameter of $3.8$ meters. What is its radius?
\end{questionText}
\choiceA{$7.6$ m}
\choiceB{$3.8$ m}
\choiceC{$1.9$ m}
\choiceD{$0.95$ m}
\correctAnswer{C}
\explanation{The radius is half the diameter: $r = d \div 2 = 3.8 \div 2 = 1.9$ meters.}
\end{practiceQuestion}

\begin{practiceQuestion}{7-1-q10}{mc}
\begin{questionText}
Which of the following is NOT a part of a circle?
\end{questionText}
\choiceA{Radius}
\choiceB{Chord}
\choiceC{Vertex}
\choiceD{Diameter}
\correctAnswer{C}
\explanation{Radius, chord, and diameter are all parts of a circle. A vertex is a corner of a polygon, not a part of a circle.}
\end{practiceQuestion}

\begin{practiceQuestion}{7-1-q11}{mc}
\begin{questionText}
A circular fountain has a diameter of $40$ feet. A person walks from the edge of the fountain straight to the center. How far does the person walk?
\end{questionText}
\choiceA{$10$ feet}
\choiceB{$20$ feet}
\choiceC{$40$ feet}
\choiceD{$80$ feet}
\correctAnswer{B}
\explanation{The distance from the edge to the center is the radius. The radius is $r = d \div 2 = 40 \div 2 = 20$ feet.}
\end{practiceQuestion}

\begin{practiceQuestion}{7-1-q12}{mc}
\begin{questionText}
If the radius of a circle is tripled, what happens to its diameter?
\end{questionText}
\choiceA{It stays the same.}
\choiceB{It doubles.}
\choiceC{It triples.}
\choiceD{It becomes six times as large.}
\correctAnswer{C}
\explanation{Since $d = 2r$, the diameter is always twice the radius. If the radius triples to $3r$, the new diameter is $2 \times 3r = 6r = 3(2r) = 3d$, so the diameter also triples.}
\end{practiceQuestion}

\begin{practiceQuestion}{7-1-q13}{mc}
\begin{questionText}
A decorating ribbon is wrapped once around the outside of a circular cake. The ribbon represents which part of the circle?
\end{questionText}
\choiceA{Radius}
\choiceB{Diameter}
\choiceC{Chord}
\choiceD{Circumference}
\correctAnswer{D}
\explanation{The circumference is the distance around a circle. A ribbon wrapped once around the outside of the cake represents the circumference.}
\end{practiceQuestion}

\begin{practiceQuestion}{7-1-q14}{mc}
\begin{questionText}
Which statement about a diameter is true?
\end{questionText}
\choiceA{A diameter is shorter than any chord.}
\choiceB{A diameter is the longest chord of a circle.}
\choiceC{A diameter does not pass through the center.}
\choiceD{A diameter equals half the radius.}
\correctAnswer{B}
\explanation{A diameter is a chord that passes through the center of the circle, and it is always the longest possible chord in that circle.}
\end{practiceQuestion}

\begin{practiceQuestion}{7-1-q15}{mc}
\begin{questionText}
A circle has a radius of $\frac{7}{2}$ inches. What is the diameter?
\end{questionText}
\choiceA{$\frac{7}{4}$ inches}
\choiceB{$\frac{7}{2}$ inches}
\choiceC{$7$ inches}
\choiceD{$14$ inches}
\correctAnswer{C}
\explanation{The diameter is twice the radius: $d = 2 \times \frac{7}{2} = 7$ inches.}
\end{practiceQuestion}

% --- Short Answer Questions (6) ---

\begin{practiceQuestion}{7-1-q16}{sa}
\begin{questionText}
A circular pond has a diameter of $34$ feet. What is the radius of the pond?
\end{questionText}
\correctAnswer{$17$ feet}
\explanation{The radius is half the diameter: $r = d \div 2 = 34 \div 2 = 17$ feet.}
\end{practiceQuestion}

\begin{practiceQuestion}{7-1-q17}{sa}
\begin{questionText}
The radius of a circular garden is $12.5$ meters. What is the diameter?
\end{questionText}
\correctAnswer{$25$ meters}
\explanation{The diameter is twice the radius: $d = 2r = 2 \times 12.5 = 25$ meters.}
\end{practiceQuestion}

\begin{practiceQuestion}{7-1-q18}{sa}
\begin{questionText}
A semicircle has a straight edge (flat side) of $28$ centimeters. What is the radius of the full circle?
\end{questionText}
\correctAnswer{$14$ cm}
\explanation{The straight edge of a semicircle is the diameter of the full circle. The radius is half the diameter: $r = 28 \div 2 = 14$ cm.}
\end{practiceQuestion}

\begin{practiceQuestion}{7-1-q19}{sa}
\begin{questionText}
A circular table has a diameter of $5$ feet. What is the radius of the table in inches? (1 foot $= 12$ inches)
\end{questionText}
\correctAnswer{$30$ inches}
\explanation{Convert the diameter to inches: $5 \times 12 = 60$ inches. The radius is half the diameter: $r = 60 \div 2 = 30$ inches.}
\end{practiceQuestion}

\begin{practiceQuestion}{7-1-q20}{sa}
\begin{questionText}
The diameter of a round window is $0.6$ meters. What is the radius of the window in centimeters? ($1$ m $= 100$ cm)
\end{questionText}
\correctAnswer{$30$ cm}
\explanation{Convert the diameter to centimeters: $0.6 \times 100 = 60$ cm. The radius is half the diameter: $r = 60 \div 2 = 30$ cm.}
\end{practiceQuestion}

\begin{practiceQuestion}{7-1-q21}{sa}
\begin{questionText}
Circle A has a radius of $8$ cm and Circle B has a radius of $5$ cm. How much longer is the diameter of Circle A than the diameter of Circle B?
\end{questionText}
\correctAnswer{$6$ cm}
\explanation{Diameter A $= 2 \times 8 = 16$ cm. Diameter B $= 2 \times 5 = 10$ cm. The difference is $16 - 10 = 6$ cm.}
\end{practiceQuestion}

% --- Graphical Questions (3 GMC + 3 GSA) ---

\begin{practiceQuestion}{7-1-q22}{gmc}
\begin{questionText}
Look at the circle below. Which labeled segment is a radius?

\begin{center}
\begin{tikzpicture}
  \draw[thick] (0,0) circle (2);
  \fill (0,0) circle (2pt) node[below left] {$O$};
  % Radius OA
  \draw[thick,funBlue] (0,0) -- (2,0);
  \fill (2,0) circle (2pt) node[right] {$A$};
  \node[funBlue,below] at (1,0) {Segment 1};
  % Chord BC (not through center)
  \fill (-1.4,1.4) circle (2pt) node[above left] {$B$};
  \fill (1.4,1.4) circle (2pt) node[above right] {$C$};
  \draw[thick,funOrange] (-1.4,1.4) -- (1.4,1.4);
  \node[funOrange,above] at (0,1.4) {Segment 2};
  % Diameter DE
  \fill (-2,0) circle (2pt) node[left] {$D$};
  \fill (2,0) circle (2pt);
  \draw[thick,funGreen] (-2,0) -- (2,0);
  \node[funGreen,above] at (-1,0) {Segment 3};
\end{tikzpicture}
\end{center}
\end{questionText}
\choiceA{Segment 1 only}
\choiceB{Segment 2 only}
\choiceC{Segment 3 only}
\choiceD{Segments 1 and 2}
\correctAnswer{A}
\explanation{A radius connects the center of the circle to a point on the circle. Segment 1 ($OA$) runs from center $O$ to point $A$ on the circle, so it is the radius. Segment 2 is a chord and Segment 3 is a diameter.}
\end{practiceQuestion}

\begin{practiceQuestion}{7-1-q23}{gmc}
\begin{questionText}
In the circle below, which segment is a chord that is NOT a diameter?

\begin{center}
\begin{tikzpicture}
  \draw[thick] (0,0) circle (2);
  \fill (0,0) circle (2pt) node[below right] {$O$};
  % Diameter PQ
  \fill (0,2) circle (2pt) node[above] {$P$};
  \fill (0,-2) circle (2pt) node[below] {$Q$};
  \draw[thick,funGreen] (0,2) -- (0,-2);
  \node[funGreen,right] at (0.15,0) {$\overline{PQ}$};
  % Chord RS (not through center)
  \fill (-1.73,-1) circle (2pt) node[below left] {$R$};
  \fill (1.73,-1) circle (2pt) node[below right] {$S$};
  \draw[thick,funOrange] (-1.73,-1) -- (1.73,-1);
  \node[funOrange,below] at (0,-1) {$\overline{RS}$};
  % Radius OT
  \fill (1.41,1.41) circle (2pt) node[above right] {$T$};
  \draw[thick,funBlue] (0,0) -- (1.41,1.41);
  \node[funBlue,above left] at (0.7,0.7) {$\overline{OT}$};
\end{tikzpicture}
\end{center}
\end{questionText}
\choiceA{$\overline{PQ}$}
\choiceB{$\overline{RS}$}
\choiceC{$\overline{OT}$}
\choiceD{Both $\overline{PQ}$ and $\overline{RS}$}
\correctAnswer{B}
\explanation{A chord has both endpoints on the circle. $\overline{RS}$ is a chord that does not pass through center $O$, so it is not a diameter. $\overline{PQ}$ passes through the center, making it a diameter. $\overline{OT}$ starts at the center, making it a radius.}
\end{practiceQuestion}

\begin{practiceQuestion}{7-1-q24}{gmc}
\begin{questionText}
The circle below has a radius labeled. What is the diameter of the circle?

\begin{center}
\begin{tikzpicture}
  \draw[thick] (0,0) circle (2);
  \fill (0,0) circle (2pt) node[below left] {$O$};
  \draw[thick,funBlue] (0,0) -- (2,0) node[midway,below] {$r = 7$ cm};
  \fill (2,0) circle (2pt) node[right] {$A$};
\end{tikzpicture}
\end{center}
\end{questionText}
\choiceA{$3.5$ cm}
\choiceB{$7$ cm}
\choiceC{$14$ cm}
\choiceD{$21$ cm}
\correctAnswer{C}
\explanation{The diameter is twice the radius: $d = 2r = 2 \times 7 = 14$ cm.}
\end{practiceQuestion}

\begin{practiceQuestion}{7-1-q25}{gsa}
\begin{questionText}
Look at the circle below. The diameter is labeled. What is the radius?

\begin{center}
\begin{tikzpicture}
  \draw[thick] (0,0) circle (2);
  \fill (0,0) circle (2pt) node[below] {$O$};
  \draw[thick,funBlue] (-2,0) -- (2,0) node[midway,above] {$d = 16$ cm};
  \fill (-2,0) circle (2pt) node[left] {$M$};
  \fill (2,0) circle (2pt) node[right] {$N$};
\end{tikzpicture}
\end{center}
\end{questionText}
\correctAnswer{$8$ cm}
\explanation{The radius is half the diameter: $r = d \div 2 = 16 \div 2 = 8$ cm.}
\end{practiceQuestion}

\begin{practiceQuestion}{7-1-q26}{gsa}
\begin{questionText}
The circle below has a radius labeled. What is the diameter?

\begin{center}
\begin{tikzpicture}
  \draw[thick] (0,0) circle (2.2);
  \fill (0,0) circle (2pt) node[below left] {$O$};
  \draw[thick,funBlue] (0,0) -- (0,2.2) node[midway,right] {$r = 11$ in};
  \fill (0,2.2) circle (2pt) node[above] {$K$};
\end{tikzpicture}
\end{center}
\end{questionText}
\correctAnswer{$22$ inches}
\explanation{The diameter is twice the radius: $d = 2r = 2 \times 11 = 22$ inches.}
\end{practiceQuestion}

\begin{practiceQuestion}{7-1-q27}{gsa}
\begin{questionText}
Two concentric circles (circles that share the same center) are shown below. How much longer is the diameter of the larger circle than the diameter of the smaller circle?

\begin{center}
\begin{tikzpicture}
  \draw[thick,funBlue] (0,0) circle (2.5);
  \draw[thick,funOrange] (0,0) circle (1);
  \fill (0,0) circle (2pt) node[below left] {$O$};
  \draw[thick,funBlue] (0,0) -- (2.5,0) node[midway,above] {$10$};
  \draw[thick,funOrange] (0,0) -- (-1,0) node[midway,above] {$4$};
\end{tikzpicture}
\end{center}
\end{questionText}
\correctAnswer{$12$ units}
\explanation{Large-circle diameter: $2 \times 10 = 20$. Small-circle diameter: $2 \times 4 = 8$. The difference is $20 - 8 = 12$ units.}
\end{practiceQuestion}
